\chapter{Discusión y conclusiones}
\label{cap:conclusiones}

\section{Resumen de hipótesis}
\pendiente{tabla con la evidencia de cada hipótesis (salida \texttt{resultados/resumen\_hipotesis.csv}).}

\section{Discusión}
\pendiente{interpretación de resultados; comparación con la literatura (p.~ej.\ \cite{shoeb2010}); relevancia clínica del umbral de 1 falsa alarma por hora.}

\section{Limitaciones}
\begin{itemize}
  \item Solo se usó una muestra de los registros sin crisis, lo que aumenta artificialmente la prevalencia respecto de un monitoreo continuo.
  \item La base es pediátrica y de un solo centro.
  \item El modelo clasifica ventanas independientes, sin aprovechar la continuidad temporal.
\end{itemize}

\section{Trabajo futuro}
\begin{itemize}
  \item Redes convolucionales sobre el espectrograma, tratado como imagen \parencite{gonzalez2018}.
  \item Clasificación multiclase (preictal, ictal, interictal) con salida softmax en el simplex.
  \item Suavizado temporal de las decisiones y evaluación en registros continuos completos.
\end{itemize}
